% BHU PYQ -- Auto-generated & error-corrected
\documentclass[11pt,a4paper]{article}

\usepackage[a4paper,top=2.5cm,bottom=2.5cm,left=2.8cm,right=2.8cm,headheight=14pt]{geometry}
\usepackage{amsmath,amssymb}
\usepackage[T1]{fontenc}
\usepackage[utf8]{inputenc}
\usepackage{lmodern}
\usepackage{microtype}
\usepackage{fancyhdr}
\usepackage{enumitem}
\usepackage{hyperref}
\usepackage{lastpage}
\usepackage{parskip}

\hypersetup{colorlinks=true,urlcolor=black,linkcolor=black,
  pdftitle={STA-301: Statistical Inference -- BHU PYQ},pdfauthor={Banaras Hindu University}}

\pagestyle{fancy}\fancyhf{}
\renewcommand{\headrulewidth}{0.4pt}\renewcommand{\footrulewidth}{0.4pt}
\fancyhead[L]{\small   \textit{Banaras Hindu University}}
\fancyhead[R]{\small   \textit{STA-301: Statistical Inference}}
\fancyfoot[C]{\small Page  \thepage\ of \pageref{LastPage}}


\newlist{parts}{enumerate}{2}
\setlist[parts,1]{label=(\alph*),leftmargin=*,itemsep=5pt,topsep=3pt}
\setlist[parts,2]{label=(\roman*),leftmargin=*,itemsep=3pt,topsep=2pt}

\newcommand{\pts}[1]{\hfill   \textit{[\#1]}}


\begin{document}


\begin{center}
  {\large\bfseries BANARAS HINDU UNIVERSITY}\\[2pt]
  {\normalsize B.A. (Hons.) Semester III Examination, 2023-24}\\[6pt]
  \rule{0.8\linewidth}{0.5pt}\\[6pt]
  {\large\bfseries Paper No. STA-301: Statistical Inference}\\[4pt]
  {\normalsize Time: 3 Hours\hfill Maximum Marks: 70}
\end{center}
\rule{\linewidth}{0.4pt}
\smallskip



\noindent   \textit{Instructions: Answer   \textbf{five} questions including Question~1 which is compulsory. Symbols have their usual meanings.}
\bigskip

\medskip


\noindent   \textbf{Question 1.}

\begin{parts}
  \item Explain the concept of the sampling distribution of a statistic and standard error. Describe the utility of standard error in hypothesis testing. Derive the expressions for the standard error of the following cases:
  
\begin{parts}
    \item The mean of a random sample of size $n$ obtained from a normal distribution.
    \item The difference of the means of two independent random samples of sizes $n_1$ and $n_2$ from normal distributions $N(\mu_1, \sigma^2)$ and $N(\mu_2, \sigma^2)$ respectively.
  \end{parts}
  Also discuss the case of equal but unknown variance for both populations. \pts{8}
  \item Obtain the sampling distribution of the sum of independent Poisson random variates. \pts{6}
\end{parts}
\medskip\rule{\linewidth}{0.2pt}

\medskip


\noindent   \textbf{Question 2.}

\begin{parts}
  \item Discuss the criteria of a good estimator. Let $X_1, X_2, X_3, X_4$ be a random sample from a $N(\mu, \sigma^2)$ population. Find the variance of $T = \frac{X_1 + 3X_2 + 2X_3 + X_4}{7}$ and $ar{X} = \frac{X_1 + X_2 + X_3 + X_4}{4}$. Which one is preferred between $T$ and $ar{X}$ for the estimation of the mean of a normal distribution and why? \pts{7}
  \item Show that for a normal distribution, the sample variance $S^2 = \frac{1}{n}\sum_{i=1}^n (X_i - ar{X})^2$ is not an unbiased estimate of the population variance $\sigma^2$. Obtain an unbiased estimate of the population variance. \pts{7}
\end{parts}
\medskip\rule{\linewidth}{0.2pt}

\medskip


\noindent   \textbf{Question 3.}

\begin{parts}
  \item Let $P_1$ and $P_2$ be the unknown proportions of students wearing glasses in two universities $A$ and $B$. To compare $P_1$ and $P_2$, samples of sizes $n_1$ and $n_2$ are taken from the two populations and the number of students wearing glasses is found to be $x_1$ and $x_2$ respectively. Suggest an unbiased estimate of $P_1 - P_2$ and obtain its sampling distribution when $n_1$ and $n_2$ are large. Hence, explain how to test the hypothesis $H_0 : P_1 = P_2$. \pts{7}
  \item Twenty people were attacked by a disease and only 18 survived. Will you reject the hypothesis that the survival rate, if attacked by this disease, is 85\% in favor of the hypothesis that it is more, at 5\% level of significance? (Use large sample test; the tabulated $Z$ value in a one-tailed test is 1.645). \pts{7}
\end{parts}
\medskip\rule{\linewidth}{0.2pt}

\medskip


\noindent   \textbf{Question 4.}



\noindent Describe the $\chi^2$ test of significance and state its various applications. What are the assumptions of the $\chi^2$ test of goodness of fit? Discuss its additive property and the limiting form of $\chi^2$ for large degrees of freedom. \pts{14}
\medskip\rule{\linewidth}{0.2pt}

\medskip


\noindent   \textbf{Question 5.}

\begin{parts}
  \item Explain stating clearly the assumptions involved in the $t$-test for testing the significance of the difference between two means. Find the value of the $t$-statistic in a sample of size eight: $-4$, $-2$, $-2$, $0$, $2$, $2$, $3$, $3$, taking the mean of the universe to be zero. How would you proceed further to test at 5\% level of significance? (Given $t_{0.05, 7} = 2.365$). \pts{7}
  \item Discuss the $F$-distribution along with its applications. \pts{7}
\end{parts}
\medskip\rule{\linewidth}{0.2pt}

\medskip


\noindent   \textbf{Question 6.}

\begin{parts}
  \item State the Weak Law of Large Numbers (WLLN). Examine whether the WLLN holds or not for the sequence $\{X_n\}$ of independent random variables defined as follows:
  \[
  P(X_n = 2^n) = 2^{-2n-1}, \quad P(X_n = -2^n) = 2^{-2n-1}, \quad P(X_n = 0) = 1 - 2^{-2n}
  \] \pts{7}
  \item Write the statement of the Central Limit Theorem (CLT) for the i.i.d. case. Also explain its utility. \pts{7}
\end{parts}
\medskip\rule{\linewidth}{0.2pt}

\medskip


\noindent   \textbf{Question 7.}

\begin{parts}
  \item Define order statistics. What are the cumulative distribution functions of the smallest and largest order statistics? Also, obtain the probability density function of the $r$-th order statistic. \pts{7}
  \item Discuss the applications of non-parametric tests, along with their advantages and disadvantages. Also, explain the Run test. \pts{7}
\end{parts}
\medskip\rule{\linewidth}{0.2pt}

\medskip


\noindent   \textbf{Question 8.}



\noindent Write short notes on any   \textbf{three} of the following: \pts{14}

\begin{parts}
  \item Errors, Critical Region, Level of Significance, and Power of a test in hypothesis testing.
  \item $\chi^2$-test for independence of attributes.
  \item Kolmogorov-Smirnov (K-S) test.
  \item Fisher's $Z$-transformation and its applications.
\end{parts}

\vfill
\rule{\linewidth}{0.4pt}

\begin{center}\small
  \href{https://bhu-exam.vercel.app/}{Click here to visit our website}\\[4pt]\href{https://me-aryan.vercel.app/}{Click here to visit our contributor}\\[4pt]\href{https://quashan-me.vercel.app/}{Click here to visit our contributor}
\end{center}

\end{document}